\documentclass[12pt]{article}

% Idioma e codificação
\usepackage[brazil]{babel}
\usepackage[utf8]{inputenc}
\usepackage[T1]{fontenc}

% Matemática, listas e tabelas
\usepackage{amsmath, amssymb}
\usepackage{enumitem}
\usepackage{float} 
\usepackage{tabularx}
\usepackage{booktabs}
\usepackage{colortbl}

% Layout e cores
\usepackage{geometry}
\usepackage{graphicx}
\usepackage{xcolor}

% Cores institucionais
\definecolor{maincolor}{RGB}{0,100,50}
\definecolor{forestgreen}{RGB}{0,100,50}
\definecolor{tableheader}{RGB}{230,240,230}

% Cabeçalho e rodapé
\usepackage{fancyhdr}

% Blocos
\usepackage[most]{tcolorbox}
\newtcolorbox{block}[1]{%
  colback=green!5,
  colframe=maincolor,
  coltitle=white,
  title=#1,
  fonttitle=\bfseries,
  boxrule=1.5pt,
  arc=3mm,
  left=6pt,
  right=6pt,
  top=6pt,
  bottom=6pt
}

% Margens
\geometry{margin=2.5cm}

% Fancyhdr
\pagestyle{fancy}
\fancyhf{}
\renewcommand{\headrulewidth}{2pt}
\renewcommand{\footrulewidth}{2pt}
\renewcommand{\headrule}{\hbox to\headwidth{\color{forestgreen}\leaders\hrule height \headrulewidth\hfill}}
\renewcommand{\footrule}{\hbox to\headwidth{\color{forestgreen}\leaders\hrule height \footrulewidth\hfill}}
\setlength{\headheight}{52pt}
\addtolength{\topmargin}{-20pt}

\fancyhead[L]{%
  \includegraphics[height=1.2cm]{Universidade_de_Rio_Verde.jpg}
}

\fancyhead[R]{%
  \raggedleft
  \textcolor{forestgreen}{%
    \small
    \textbf{Lista de Exercícios}\\
    Me.~Sergio Souza Novak\\
    \textit{Banco de Dados}
    }
}
    
\fancyfoot[C]{\textcolor{forestgreen}{\thepage}}
    
\begin{document}
    
\vspace*{0.2cm}

\section*{Exercício 01}

Analise a instância do banco de dados abaixo para o sistema \textbf{PetShop 2.0}. As tabelas demonstram como os dados estão organizados e relacionados.

\begin{table}[H]
    \centering
    \small
    % Tabela Pessoa
    \textbf{PESSOA}\\
    \begin{tabular}{|l|l|l|}
    \hline
    \rowcolor{tableheader} \underline{cpf} (PK) & nome & email \\ \hline
    111.222.333-44 & João Silva & joao@email.com \\ \hline
    555.666.777-88 & Ana Costa & ana@email.com \\ \hline
    999.000.111-22 & Ricardo Souza & ricardo@email.com \\ \hline
    \end{tabular}

    \vspace{0.5cm}

    % Tabela Veterinário e Gerente
    \begin{minipage}{0.45\textwidth}
        \centering
        \textbf{VETERINARIO}\\
        \begin{tabular}{|l|l|}
        \hline
        \rowcolor{tableheader} \underline{\textit{cpf\_vet}} (FK) & crmv \\ \hline
        111.222.333-44 & MG-1234 \\ \hline
        \end{tabular}
    \end{minipage}
    \begin{minipage}{0.45\textwidth}
        \centering
        \textbf{GERENTE}\\
        \begin{tabular}{|l|l|}
        \hline
        \rowcolor{tableheader} \underline{\textit{cpf\_gerente}} (FK) & \textit{cpf\_sup} (FK) \\ \hline
        555.666.777-88 & NULL \\ \hline
        999.000.111-22 & 555.666.777-88 \\ \hline
        \end{tabular}
    \end{minipage}

    \vspace{0.5cm}

    % Tabela Animal e Vacina
    \textbf{ANIMAL}\\
    \begin{tabular}{|l|l|l|}
    \hline
    \rowcolor{tableheader} \underline{id\_animal} (PK) & nome\_pet & \textit{cpf\_cliente} (FK) \\ \hline
    10 & Rex (Cão) & 111.222.333-44 \\ \hline
    20 & Mingau (Gato) & 111.222.333-44 \\ \hline
    \end{tabular}

    \vspace{0.5cm}

    \textbf{HIST\_VACINA}\\
    \begin{tabular}{|l|l|l|}
    \hline
    \rowcolor{tableheader} \underline{\textit{id\_animal}} (FK) & \underline{data\_vac} & lote \\ \hline
    10 & 10/01/2026 & V8-456 \\ \hline
    10 & 10/02/2026 & RV-789 \\ \hline
    20 & 15/01/2026 & V5-123 \\ \hline
    \end{tabular}
\end{table}



\section*{Instruções do Exercício}

Com base na visualização das tabelas acima, responda:

\begin{enumerate}[label=\alph*)]
    \item Na tabela \texttt{HIST\_VACINA}, o que aconteceria se tentássemos inserir uma nova linha com o \texttt{id\_animal = 30}? Justifique com base no conceito de \textbf{Integridade Referencial} visto nos slides.
    \item Observe a tabela \texttt{GERENTE}. Explique como o valor presente na coluna \texttt{cpf\_sup} de Ricardo Souza se relaciona com a mesma tabela e o que isso representa na hierarquia da empresa.
    \item Por que o CPF de João Silva aparece tanto na tabela \texttt{PESSOA} quanto na tabela \texttt{VETERINARIO}? Qual estratégia de mapeamento de herança isso demonstra?
    \item Identifique um exemplo de relacionamento \textbf{1:N} entre as tabelas acima e explique como a "Chave Estrangeira" é usada para manter esse vínculo.
\end{enumerate}

\begin{flushright}
\textit{Use os exemplos das páginas 38 a 49 do material como referência.}
\end{flushright}
\section*{Exercício 02}

A empresa \textbf{Pé na Tábua Autopeças} necessita organizar o seu estoque. No modelo lógico, as peças são classificadas por tipo, e o armazenamento é controlado por prateleiras. Analise as tabelas abaixo que representam uma parte do banco de dados:

\begin{table}[H]
    \centering
    \small
    
    \textbf{PECA}\\
    \begin{tabular}{|l|l|l|}
    \hline
    \rowcolor{tableheader} \underline{part\_number} (PK) & descricao & fabricante \\ \hline
    P-001 & Amortecedor Dianteiro & Cofap \\ \hline
    P-002 & Filtro de Óleo & Fram \\ \hline
    P-003 & Pastilha de Freio & Bosch \\ \hline
    \end{tabular}

    \vspace{0.4cm}

    % Especializações
    \begin{minipage}{0.48\textwidth}
        \centering
        \textbf{PECA\_NACIONAL}\\
        \begin{tabular}{|l|l|}
        \hline
        \rowcolor{tableheader} \underline{\textit{part\_num}} (FK) & ipi \\ \hline
        P-001 & 8\% \\ \hline
        \end{tabular}
    \end{minipage}
    \begin{minipage}{0.48\textwidth}
        \centering
        \textbf{PECA\_IMPORTADA}\\
        \begin{tabular}{|l|l|}
        \hline
        \rowcolor{tableheader} \underline{\textit{part\_num}} (FK) & taxa\_imp \\ \hline
        P-003 & 15\% \\ \hline
        \end{tabular}
    \end{minipage}

    \vspace{0.6cm}

    \textbf{PRATELEIRA}\\
    \begin{tabular}{|l|l|l|}
    \hline
    \rowcolor{tableheader} \underline{id\_prat} (PK) & setor & corredor \\ \hline
    A1 & Suspensão & 01 \\ \hline
    B5 & Frenagem & 02 \\ \hline
    \end{tabular}

    \vspace{0.6cm}

    \textbf{ESTOQUE\_ITEM (Tabela Associativa N:N)}\\
    \begin{tabular}{|l|l|l|}
    \hline
    \rowcolor{tableheader} \underline{\textit{part\_num}} (FK) & \underline{\textit{id\_prat}} (FK) & quantidade \\ \hline
    P-001 & A1 & 15 \\ \hline
    P-003 & B5 & 22 \\ \hline
    P-001 & B5 & 5 \\ \hline
    \end{tabular}
\end{table}



\section*{Instruções do Exercício}

Analise a estrutura e os dados das tabelas acima para responder aos itens:

\begin{enumerate}[label=\alph*)]
    \item A tabela \texttt{ESTOQUE\_ITEM} possui uma \textbf{Chave Primária Composta}. Identifique quais campos a formam e explique, com base nos dados apresentados, por que não poderíamos usar apenas o \texttt{part\_num} como PK nesta tabela.
    
    \item Na tabela \texttt{ESTOQUE\_ITEM}, a peça \texttt{P-001} aparece duas vezes. Isso caracteriza uma violação de integridade ou é uma característica de um relacionamento \textbf{N:N}? Justifique sua resposta.
    
    \item Caso a peça \texttt{P-002} fosse importada, em qual tabela deveríamos inserir os dados adicionais dela? Explique como a integridade referencial garante que essa peça "exista" antes de ser classificada como importada.
    
    \item Como o conceito de \textbf{Atributo de Relacionamento} é aplicado na tabela \texttt{ESTOQUE\_ITEM}? Onde está armazenada a informação de "quantidade" e por que ela não poderia estar na tabela \texttt{PECA}?
\end{enumerate}

\begin{block}{Desafio Prático}
Se decidirmos apagar a Prateleira \texttt{B5} do sistema, o que aconteceria com os registros de estoque das peças \texttt{P-001} e \texttt{P-003} se a restrição de integridade for configurada como \textbf{ON DELETE RESTRICT}?
\end{block} 

\section*{Exercício 03}

Um hospital universitário tentou implementar seu banco de dados, mas o esquema resultante apresenta graves falhas de modelagem e violações das regras de Edgar F. Codd. Analise as tabelas abaixo, que representam o estado atual do sistema de \textbf{Internações}:

\begin{table}[H]
    \centering
    \small
    
    \textbf{MEDICO}\\
    \begin{tabular}{|l|l|l|}
    \hline
    \rowcolor{tableheader} \underline{crm} (PK) & nome & especialidade \\ \hline
    CRM-123 & Dr. Arnaldo & Cardiologia \\ \hline
    CRM-123 & Dr. Arnaldo & Cardiologia \\ \hline
    CRM-999 & Dra. Beatriz & NULL \\ \hline
    \end{tabular}

    \vspace{0.4cm}

    \textbf{PACIENTE}\\
    \begin{tabular}{|l|l|l|}
    \hline
    \rowcolor{tableheader} \underline{id\_paciente} (PK) & nome & data\_nasc \\ \hline
    101 & Paciente A & 12/05/1980 \\ \hline
    NULL & Paciente B & 20/10/1995 \\ \hline
    \end{tabular}

    \vspace{0.4cm}

    \textbf{INTERNACAO (Tabela com falhas de integridade)}\\
    \begin{tabular}{|l|l|l|l|}
    \hline
    \rowcolor{tableheader} \underline{id\_int} (PK) & \textit{id\_paciente} (FK) & \textit{crm\_medico} (FK) & quarto \\ \hline
    1 & 101 & CRM-123 & 202 \\ \hline
    2 & 500 & CRM-123 & 205 \\ \hline
    3 & 101 & CRM-888 & 210 \\ \hline
    \end{tabular}
\end{table}



\section*{Instruções do Exercício}

Com base nos conceitos de integridade discutidos na \textbf{Aula 04}, identifique e explique as falhas encontradas:

\begin{enumerate}[label=\alph*)]
    \item \textbf{Violação de Integridade de Entidade:} Observe a tabela \texttt{PACIENTE}. Existe um erro grave na segunda linha. Qual é esse erro e por que o Modelo Relacional proíbe tal ocorrência?
    
    \item \textbf{Violação de Integridade de Domínio (Vazio):} Na tabela \texttt{MEDICO}, a Dra. Beatriz está com a especialidade como \texttt{NULL}. Considerando que todo médico deve obrigatoriamente ter uma especialidade no sistema, como essa restrição deveria ter sido definida no projeto lógico?
    
    \item \textbf{Violação de Integridade Referencial (O erro mais grave):} Analise as linhas 2 e 3 da tabela \texttt{INTERNACAO}.
    \begin{itemize}
        \item O que há de errado com o \texttt{id\_paciente = 500}?
        \item O que há de errado com o \texttt{crm\_medico = CRM-888}?
    \end{itemize}
    Explique o conceito de "Chave Estrangeira Órfã" aplicado a este caso.
    
    \item \textbf{Redundância e Chaves:} Na tabela \texttt{MEDICO}, o Dr. Arnaldo aparece duas vezes com o mesmo CRM. Isso é permitido em uma coluna definida como \textit{Primary Key}? O que acontece com a unicidade da tupla?
\end{enumerate}

\begin{block}{Desafio de Correção}
Como você redesenharia a tabela \texttt{INTERNACAO} para garantir que nunca fosse possível internar um paciente que não existe ou atribuir um médico não cadastrado? Cite o uso de \textit{Constraints} (Restrições).
\end{block}


\section*{Exercício 04}

A correta alocação da Chave Estrangeira (FK) é fundamental para evitar a redundância de dados e o excesso de valores nulos (\texttt{NULL}). Com base nas definições de mapeamento lógico, determine em qual tabela a FK deve ser inserida para cada cenário abaixo.

\begin{block}{Instrução}
Para cada item, indique a \textbf{Tabela de Origem}, a \textbf{Tabela de Destino} da chave e justifique sua escolha com base na cardinalidade (\textit{1:1, 1:N ou N:N}).
\end{block}

\begin{enumerate}[label=\textbf{\arabic*.}]
    \item \textbf{Cenário de Biblioteca:} Uma biblioteca possui \textbf{Livros} e \textbf{Editoras}. Uma editora pode publicar vários livros, mas cada livro é publicado por apenas uma única editora. 
    \begin{itemize}
        \item Onde fica a FK? \hrulefill
    \end{itemize}

    \item \textbf{Cenário de Gestão Pública:} Um \textbf{Prefeito} governa exatamente um \textbf{Município}, e cada município possui apenas um prefeito em exercício. Sabe-se que todos os municípios cadastrados possuem um prefeito, mas nem todo político cadastrado como "Prefeito" está atualmente governando uma cidade.
    \begin{itemize}
        \item Onde fica a FK para evitar valores nulos? \hrulefill
    \end{itemize}

    \item \textbf{Cenário de E-commerce:} Um \textbf{Cliente} realiza vários \textbf{Pedidos}, e cada pedido pertence a um único cliente.
    \begin{itemize}
        \item Onde fica a FK? \hrulefill
    \end{itemize}

    \item \textbf{Cenário Escolar:} Um \textbf{Aluno} pode se matricular em várias \textbf{Disciplinas}, e uma disciplina pode ter vários alunos matriculados.
    \begin{itemize}
        \item Como este relacionamento é resolvido no modelo lógico? \hrulefill
    \end{itemize}

    \item \textbf{Cenário de Recursos Humanos:} Em uma empresa, um \textbf{Departamento} possui um \textbf{Gerente} (que é um empregado). Cada departamento tem apenas um gerente, e cada gerente gerencia apenas um departamento. Sabe-se que todos os departamentos devem ter um gerente obrigatoriamente.
    \begin{itemize}
        \item Em qual tabela a FK deve ser propagada? \hrulefill
    \end{itemize}
\end{enumerate}


\vfill
\begin{center}
    \textit{"A chave estrangeira migra para o lado N da relação ou para o lado da participação total no caso 1:1."}
\end{center}
\end{document}